\section{Sparse discovery, independent replay, and the old planner}
\label{sec:algorithms}

\subsection{An immutable-source algorithm}

Peeling needs no evolving matrix. For every source column $h$, maintain the
number $d_h$ of incident live rows and the identity of its sole owner when
$d_h=1$. An incident row with entry at least two does not make a unit leaf.
When row $r$ is removed, visit its source nonzero columns once, decrease
their incident-row counts, and enqueue newly exposed unit leaves.

\begin{lstlisting}[caption={Complete raw cancellation from source incidence.}]
build row-to-column incidence lists
for each column h:
    degree[h] = number of nonzero incident rows
    owner[h] = XOR of those row labels
enqueue h if degree[h] == 1 and N[owner[h], h] == 1

while the queue is nonempty:
    pop h
    if degree[h] != 1: continue
    r = owner[h]
    record the cancellation (r,h)
    mark r deleted
    for each nonzero source column c of row r:
        degree[c] -= 1
        owner[c] ^= r
        if degree[c] == 1 and N[owner[c],c] == 1:
            enqueue c

return the schedule and the remaining row set
\end{lstlisting}

XOR is only an efficient owner-recovery implementation. Dense row labels
make it a word operation; an incident-list representation gives the same
mathematical algorithm. An all-zero row is retained. A column that is
initially absent is irrelevant. The queue may contain stale entries, which
are ignored after their degrees cease to be one.

\begin{theorem}[Sparse structural bound]\label{thm:linear}
After indexing the rows and columns, the complete raw-cancellation decision,
a successful schedule or a canonical stopping set, and the associated
sparse source data require $O(I+m+n)$ structural operations and
$O(I+m+n)$ words of storage. A successful schedule has at most $m$ pairs.
If labels use $O(\log(m+n+2))$ bits and the binary input length is $\Lambda$,
a deterministic dense-array implementation has bit cost
\[
 O\bigl(\Lambda+(I+m+n)\log(m+n+2)\bigr)
\]
under elementary word arithmetic. Sorting arbitrary external labels may
add a logarithmic preprocessing factor.
\end{theorem}

\begin{proof}
Each nonzero is read once during construction and visited at most once
when its row is deleted. A column's live degree decreases monotonically,
so it becomes degree one at most once. It is consequently enqueued at
most once, including initial leaves. Queue processing is linear in the
number of appearing columns. All other work is proportional to the row
and column arrays or the schedule. Entries are capped at two while reading;
the later algorithm does not perform arithmetic on their original magnitudes.
Counters and row labels require logarithmic space. These observations give
both structural and bit bounds.
\end{proof}

The provided core producer follows this sparse queue design. The standalone
saturated-bundle producer calls the same public sparse planner. Python dictionaries
are implementation data structures, not a proof of adversarial constant-time
hashing. The deterministic bound above describes the indexed algorithm;
the empirical measurements concern the supplied Python implementation.

\subsection{Positive certificates without producer replay}

Let a proposed schedule be
$((r_0,g_0),\ldots,(r_{s-1},g_{s-1}))$. Define $\tau(r_i)=i$ and
$\tau(r)=+\infty$ for every unlisted row. No row or pivot column may repeat.

\begin{theorem}[Timestamp certificate]\label{thm:timestamp}
The schedule is a valid no-fill column-leaf prefix if and only if, for each
$i$, the source entry $n_{r_i g_i}$ equals one and every different row $r$
with $n_{rg_i}>0$ satisfies $\tau(r)<i$. It is a complete certificate if,
in addition, every source row is listed. These conditions can be checked
by one scan of the source incidence and the schedule.
\end{theorem}

\begin{proof}
At step $i$, the live source rows are exactly those with deletion time
at least $i$. The stated inequality therefore says that column $g_i$
has only row $r_i$ left and that its entry is one. Previous column-leaf
pivots have created no fill, so the source incidence on live rows is the
actual incidence. Induction proves validity. Conversely, each genuine
leaf step forces those source incidences to belong to earlier deleted
rows. The extra coverage condition is exactly complete cancellation.
\end{proof}

The independent checker does not call the producer, maintain its queue,
or carry out \eqref{eq:positive-pivot}. It builds row timestamps and checks
source inequalities. The certificate occupies
$O(m\log(m+n+2))$ bits, apart from externally supplied source identifiers.
The same checker can verify a proper prefix by retaining $+\infty$ on
unlisted rows.

\subsection{Obstruction certificates}

An obstruction certificate supplies a nonempty, duplicate-free list $R$
of source row labels. The checker accumulates
$\min(2,\sum_{r\in R}n_{rh})$ for each appearing column and rejects if
any result is one. Theorem~\ref{thm:persistence} proves that an accepted
certificate rules out all full positive raw schedules for those input
counts. The checker need not trust the producer's claim that $R$ is
canonical; any nonempty stopping set suffices.

This distinction is useful for independent verification. Canonical discovery
is a deterministic algorithmic convenience, whereas the obstruction theorem
is a short local certificate. Relabelling the source changes labels, not
the existence of the certified stopping set. Changing the disk system can
change the matrix and may remove its obstruction.

\subsection{Why the retained capped planner was already complete}

The source implementation's legacy capped planner scores a candidate
$n_{rg}=1$ by
\[
 \left(\sum_h n_{rh}-1\right)
 \left(\sum_i n_{ig}-1\right)
\]
using capped entries. A zero score means that the donor is a unit row
leaf or that the pivot column is a unit column leaf. In either case no
surviving entry receives fill. The planner chooses a minimum score.

\begin{proposition}[Completeness of zero-fill priority]
\label{prop:old-complete}
A producer that always chooses a zero-fill raw pivot when one exists is
complete for full positive row cancellation. In particular, the retained
min-fill capped planner already had this completeness property, despite
its earlier documentation.
\end{proposition}

\begin{proof}
If the current matrix is fully cancellable, Theorem~\ref{thm:complete}
supplies a unit column leaf, so a zero-fill pivot exists. A column-leaf
pivot preserves full cancellability: any stopping set of its reduced
matrix would already be a stopping set of the original matrix, since the
removed pivot column is zero on the remaining rows.

Consider instead a unit row leaf $r$ at column $g$. If the reduced matrix
had a stopping set $R$, examine its source sum at $g$. If that sum is zero
or at least two, $R$ itself obstructs the original matrix. If it is one,
then $R\cup\{r\}$ obstructs it: the sum at $g$ becomes two and the row
leaf contributes nothing elsewhere. Either case contradicts full
cancellability. Hence this zero-fill pivot is safe as well.

Induction shows that zero-fill priority succeeds whenever any positive
schedule succeeds. Its successful outputs are always valid schedules, so
the converse implication requires no further argument.
\end{proof}

This proof does not assert that \emph{every} arbitrary singleton choice
preserves full cancellability. It establishes precisely the property used
by the retained producer. The native change removes repeated candidate
scans and unneeded positive updates, while preserving its complete-success
set. Legacy certificates remain independently replayable.

\subsection{A family exposing the avoidable scan cost}

For $m\ge2$, consider the $m\times(m+1)$ bidiagonal chain with
$n_{r,r}=n_{r,r+1}=1$ and all other entries zero:
\[
 N_m=\begin{pmatrix}
 1&1&0&\cdots&0\\
 0&1&1&\cdots&0\\
 \vdots&\ddots&\ddots&\ddots&\vdots\\
 0&\cdots&0&1&1
 \end{pmatrix}.
\]
It has $I=2m$ source incidences, no initial unit row leaf, and exposed unit
columns at the ends. Peeling successively exposes the next column.
A planner that rescans every surviving incidence at each
step performs $\Theta(m^2)$ source visits. The queue algorithm performs
$\Theta(m)$ visits. Both find a complete certificate; their distinction
is avoidable repeated work.

This is a matrix family, not a benchmark of knot recognition. It has a
geometric interpretation as a primitive attaching family on a handlebody
with one remaining handle, but an arbitrary caller-provided copy of the
matrix does not authenticate that interpretation. The genuine source-cut
audit in Section~\ref{sec:measurements} supplies the separate geometric
validation for the native paths.

\subsection{An exact boundary between useful and wasted raw search}

After a verified stopping set is found, further enumeration of positive
pivot orders cannot establish complete cancellation. This is a mathematical
termination rule for the raw subproblem, not a heuristic timeout. A higher
level search must change something outside this model: the geometric disk
system, the forest, the attaching geometry, or the decomposition surface.

For this reason the residual row set is a useful search diagnostic. It names
the curves whose joint incidence has no unique crossing. It does not measure
all remaining topology, and its cardinality is not a global hierarchy
potential. A strategy that treats the residual size as a guaranteed decreasing
topological complexity would need a further theorem.
